\documentclass[../../../include/open-logic-section]{subfiles}

\begin{document}

\olfileid{sth}{ordinals}{wo}
\olsection{حُسنِ ترتیب}

بنیادی تصور ایہ اے:

\begin{defn}
تعلق $<$، $A$ نوں \emph{حُسنِ ترتیب وچ رکھدا اے} اُدوں تے صرف اُدوں، جدوں ایہ دو شرطاں پوری کرے:
\begin{enumerate}
	\item $<$ جُڑیا ہویا اے، یعنی ہر $a, b \in A$ لئی یا $a < b$، یا
	$a = b$، یا $b < a$۔
	\item $A$ دے ہر غیر خالی ذیلی سیٹ وچ $<$ دے لحاظ نال اک ایسا !!{element} جیہدے توں کوئی ہور رکن چھوٹا نہیں اے۔ یعنی جے $\emptyset \neq X \subseteq A$، تاں
	$(\exists m \in X)(\forall z \in X)z \nless m$۔
\end{enumerate}
\end{defn}

ایہ ویکھنا سوکھا اے کہ ہُنے ویکھیاں تنّے مثالاں واقعی حُسنِ ترتیب والے تعلق سن۔

\begin{prob}
\Olref[sth][ordinals][idea]{sec} وچ قدرتی عدداں دیاں تن ترتیباں مثال دے طور اُتے
دتیاں سن۔ پرکھو کہ ہر اک حُسنِ ترتیب اے۔
\end{prob}

حُسنِ ترتیب بارے کجھ بنیادی پر بہت ضروری گلاں ہِیٹھاں نیں۔

\begin{prop}\ollabel{wo:strictorder}
جے $<$، $A$ نوں حُسنِ ترتیب وچ رکھدا اے، تاں $A$ دے ہر غیر خالی ذیلی
سیٹ وچ $<$ دے لحاظ نال اکّو اک سبھ توں چھوٹا رکن اے، تے $<$ غیر انعکاسی،
غیر متقابل تے متعدی اے۔
\end{prop}

\begin{proof}
جے $X$، $A$ دا اک غیر خالی ذیلی سیٹ اے، تاں اوہدے وچ $<$ دے لحاظ نال
اک ایسا !!{element} $m$ اے جیہدے توں کوئی ہور رکن چھوٹا نہیں، یعنی $(\forall z \in X)z \nless m$۔
چونکہ $<$ جُڑیا ہویا اے، $(\forall z \in X)m \leq z$۔
سو $m$، $<$ دے لحاظ نال $X$ دا سبھ توں چھوٹا !!{element} اے۔

غیر انعکاسیت لئی $a \in A$ مقرر کرو۔ $<$ دے لحاظ نال $\{a\}$ دا سبھ توں
چھوٹا !!{element} $a$ اے، سو $a \nless a$۔ تعدّی لئی، جے $a < b < c$،
تاں چونکہ $\{a, b, c\}$ وچ $<$ دے لحاظ نال سبھ توں چھوٹا !!{element} اے،
ایس لئی $a < c$۔ غیر تقابل، غیر انعکاسیت تے تعدّی توں نکلدا اے۔
\end{proof}

\begin{prop}\ollabel{propwoinduction}
جے $<$، $A$ نوں حُسنِ ترتیب وچ رکھدا اے، تاں ہر فارمولے $\phi(x)$ لئی:
\[
	\text{ جے }(\forall a \in A)((\forall b < a)\phi(b) \lif
		\phi(a))\text{، تاں }(\forall a \in A)\phi(a).
\]
\end{prop}

\begin{proof}
اسیں عکسِ نقیض ثابت کراں گے۔ منّ لو $\lnot(\forall a \in A)\phi(a)$، یعنی
$X = \Setabs{x \in A}{\lnot\phi(x)} \neq \emptyset$۔ تاں $X$ دا
$<$ دے لحاظ نال اک ایسا !!{element}، $a$، اے جیہدے توں کوئی ہور رکن چھوٹا نہیں۔ سو
$(\forall b < a)\phi(b)$، پر $\lnot \phi(a)$۔ \end{proof}\noindent
ایس آخری خاصیت توں قدرتی عدداں اُتے مضبوط استقرا دا اصول یاد آؤنا چاہیدا اے۔
یعنی جے $(\forall n \in \omega)((\forall m < n)\phi(m) \lif \phi(n))$،
تاں $(\forall n \in \omega)\phi(n)$۔ ایہ خاصیت حُسنِ ترتیب نوں
اک بڑا \emph{مضبوط} تصور بناندی اے۔\footnote{یاد رہے کہ سارے فارمولیاں دے
پیرامیٹر ہو سکدے نیں (جدوں تک صاف طور اُتے کجھ ہور نہ آکھیا جاوے)۔}

\end{document}
